%%% ============================================================
%%% The Three-Substrate Joint-Closure Chain on Z/10Z:
%%%   Eight Shells Survive Across (TSML, BHML, CL_STD)
%%%   with Lens-Dependence Internal to TSML at Size 7
%%%
%%% Brayden R. Sanders, M. Gish
%%% 7Site LLC, Hot Springs AR, 2026
%%%
%%% Target venue:  Mathematical Intelligencer
%%% Source:        WP115 (chain count corrected 2026-05-05);
%%%                Atlas/META_PLAN_2026-05-06/SUBSTRATE_FUNCTION_MAP/
%%%                  SFM_FINDINGS_v1.md (Q1, Q6 results, 2026-05-08);
%%%                Atlas/META_PLAN_2026-05-06/TSML_RECONCILIATION.md;
%%%                Atlas/LENS_TAXONOMY_2026-05-06/TIER_CONFLATION_AUDIT.md
%%% MSC:           08A40, 08B05, 20N02
%%% ============================================================

\documentclass[11pt,reqno]{amsart}

\usepackage{amsmath,amssymb,amsthm}
\usepackage[margin=1in]{geometry}
\usepackage{hyperref}
\usepackage{enumitem}
\usepackage{array}

\newtheorem{theorem}{Theorem}[section]
\newtheorem{lemma}[theorem]{Lemma}
\newtheorem{proposition}[theorem]{Proposition}
\newtheorem{corollary}[theorem]{Corollary}

\theoremstyle{definition}
\newtheorem{definition}[theorem]{Definition}
\newtheorem{remark}[theorem]{Remark}

\newcommand{\Z}{\mathbb{Z}}
\newcommand{\TRAW}{T_{\mathrm{RAW}}}
\newcommand{\TSYM}{T_{\mathrm{SYM}}}
\newcommand{\TSML}{T}
\newcommand{\BHML}{B}
\newcommand{\CSTD}{C}

\title[Three-substrate chain on $\Z/10\Z$]
  {The Three-Substrate Joint-Closure Chain on $\Z/10\Z$:\\
   Eight Shells Survive Across $(T,\,B,\,C)$ with
   Lens-Dependence Internal to $T$ at Size $7$}

\author{Brayden R.\ Sanders}
\address{7Site LLC, Hot Springs, Arkansas, USA}
\email{brayden@7site.co}

\author{M.\ Gish}
\address{Independent Researcher, Hot Springs, Arkansas, USA}
\email{monica.gish1992@gmail.com}

\date{2026}

\keywords{finite magma, composition lattice, joint sub-magma
  closure, three-substrate framework, lens symmetrization,
  non-commutativity, $\Z/10\Z$}

\subjclass[2020]{08A40, 08B05, 20N02}

\begin{document}

\begin{abstract}
We compute the joint sub-magma closure chain on $\Z/10\Z$ under
three composition tables drawn from the same canonical bit-pattern
encoding: $T$ (``TSML''), $B$ (``BHML''), and $C$ (``CL\_STD'').
Brute-force enumeration over all $1023$ non-empty subsets shows
that under the commutative lens $\TSYM$, subsets jointly closed
under all three tables form a strict $8$-element chain at sizes
$\{1,4,5,6,7,8,9,10\}$ — the same chain established for the pair
$(\TSYM, B)$ alone. Individually, $T = \TSYM$ admits $449$ closed
sub-magmas, $B$ admits $9$, and $C$ admits $50$; of $C$'s $50$
closed sub-magmas, $49$ are also $T$-closed and exactly $8$ are
jointly closed under all three. Adding $C$ as a third substrate
therefore preserves the entire $(T, B)$ chain without introducing
or removing a single shell. The previously-recorded lens-dependence
at size $7$ (eight shells under $(\TSYM, B)$ versus seven under
$(\TRAW, B)$, forced by the single non-commutative cell
$\TRAW(9,4)=3$) is shown here to be \emph{internal to $T$}: the
size-$7$ shell $\{0,4,5,6,7,8,9\}$ is $C$-closed regardless of the
$T$-lens choice, so $C$ does not arbitrate the asymmetric cell.
Concretely, the chain count under $(\TSYM, B, C)$ is $8$ shells
and under $(\TRAW, B, C)$ remains $7$ shells: adjoining $C$
neither cures the lens-dependence nor introduces new shells. The
four-core $\{0,7,8,9\}$ and the closed-form attractor at
$\alpha=1/2$ are lens- and table-invariant on all three
substrates. The result upgrades the joint-closure picture from a
two-table observation to a three-table structural theorem, and
locates lens choice at the correct level of the architecture.
\end{abstract}

\maketitle

\section*{Scope and tier discipline}

\textbf{PROVEN.} Subsets of $\Z/10\Z$ jointly closed under
$T$, $B$, and $C$ form an $8$-element chain at sizes
$\{1,4,5,6,7,8,9,10\}$, identical to the $(T,B)$ joint chain
(Theorems \ref{thm:three-table-chain} and
\ref{thm:lens-internal}). The four-core $\{0,7,8,9\}$ is closed
under all three tables, and the lens-dependence at size $7$
recorded in companion work is internal to the $T_{\mathrm{RAW}}$
versus $T_{\mathrm{SYM}}$ choice and does not lift to the
three-table level.

\textbf{COMPUTED.} The script \texttt{sfm\_q1\_q6\_q7.py}
(reproduced in \cite{SFM2026}) enumerates all $1023$ non-empty
subsets of $\Z/10\Z$, checks closure under each table, computes
all pairwise and three-way joint closures, and exhibits the
chain. Independently, \texttt{joint\_chain\_attractor.py}
\cite{Gen13WP115} reproduces the $(T,B)$-only counts. Total
runtime under $10$ seconds; verification deterministic.

\textbf{STRUCTURAL RHYME.} The closed-form attractor at
$\alpha=1/2$, with $H/\mathrm{Br}=1+\sqrt{3}$, lives on the
four-core (cf.~\cite{Sanders2026FourCore}) and is therefore
table-invariant for the same reason the four-core is. We do not
re-derive the attractor here; the relation $1+\sqrt{3}$ is cited
as structural motivation for caring about the four-core's
joint-closure status, not as a load-bearing identity in the
present paper.

\textbf{OPEN.} Whether other commutative composition tables on
$\Z/10\Z$ satisfying the four-core preservation criterion produce
the same $8$-shell joint chain is open. The $C$-table data point
suggests the chain is a feature of the substrate-and-four-core
structure, not of any one table; a precise statement awaits a
class-level theorem.

\textbf{Lens and substrate.} We work on $\Z/10\Z$ with the three
composition tables $T$, $B$, $C$ defined by the canonical bit
pattern of \cite{Sanders2026CLAxioms}. These choices are not
derived from first principles; they reflect a structural reading
of the substrate developed across the framework. The theorems
below are theorems on this specific $(T,B,C)$-triple. Whether
analogous tables on other substrates produce the same structure
is open. The framing follows the Drápal–Wanless (2021) line of
work \cite{DrapalWanless2021} on small finite commutative
non-associative structures.

\section{Introduction}
\label{sec:intro}

The canonical composition lattice on $\Z/10\Z$ supports three
named $10\times 10$ tables that share substrate and operator
labels but differ as multiplications. In the codebase notation,
$T$ is the TSML table (Sanders' prescribed view, $73$ HARMONY
cells), $B$ is the BHML table (Binary Hard Micro Lattice, $28$
HARMONY cells), and $C$ is the CL\_STD table (the Standard or
encoding table, $44$ HARMONY cells). The provenance and naming
are recorded in the project's foundational module
\cite{Gen13Foundations}.

This note answers the question: what is the joint sub-magma
closure structure under \emph{all three} tables simultaneously?
The answer is striking — the three-way joint chain is exactly the
two-way joint chain.

\paragraph{Statement of main theorem.}

\begin{theorem}[Three-substrate joint chain]
\label{thm:three-table-chain}
The subsets $S\subseteq\{0,\dots,9\}$ jointly closed under all
three tables $T$, $B$, $C$ form a strict $8$-element chain
\[
  \{0\} \subsetneq \{0,7,8,9\} \subsetneq \{0,6,7,8,9\}
  \subsetneq \{0,5,6,7,8,9\} \subsetneq \{0,4,5,6,7,8,9\}
  \subsetneq
\]
\[
  \{0,3,4,5,6,7,8,9\} \subsetneq \{0,2,3,4,5,6,7,8,9\}
  \subsetneq \{0,1,2,3,4,5,6,7,8,9\},
\]
of sizes $\{1,4,5,6,7,8,9,10\}$. Sizes $\{2,3\}$ are forbidden.
The chain coincides with the joint $(T,B)$ chain established in
\cite{Sanders2026FourCore}.
\end{theorem}

\paragraph{Why this is interesting.}

The three tables $T$, $B$, $C$ are independently constructed from
different structural readings of the substrate
\cite{Gen13Foundations}. Individually, $T$ has $449$ closed
sub-magmas, $B$ has $9$, and $C$ has $50$ (Lemma
\ref{lem:individual-counts}). The $T$-closures are abundant and
include many HARMONY-attracted side-branches; $B$'s few closures
include the Klein-style parity sub-magma $\{0,9\}$; $C$'s
intermediate count of $50$ records a richer encoding-respecting
closure structure. \emph{A priori} one might expect the three-way
joint closure to be much smaller than $8$, given the wildly
different individual counts. Theorem \ref{thm:three-table-chain}
says otherwise: the three-way intersection is forced to be
exactly the $(T,B)$ chain. Adding $C$ does not introduce any new
shell, and crucially does not eliminate any.

\paragraph{The lens-dependence at size $7$ is internal to $T$.}

A companion observation, recorded in
\cite{Gen13WP115,Sanders2026FourCore} and re-derived in
Section~\ref{sec:lens-internal} below, is that $T$ admits two
principled lens-symmetrization choices, $\TRAW$ (the literal
non-commutative bit pattern) and $\TSYM$ (the upper-triangle
authoritative symmetrization), differing at exactly four cells:
the asymmetric pairs $(3,9)/(9,3)$ and $(4,9)/(9,4)$.
Cell-by-cell verification yields
\[
  \TRAW(3,9) = 3,\quad \TRAW(9,3) = 7,\qquad
  \TRAW(4,9) = 7,\quad \TRAW(9,4) = 3,
\]
\[
  \TSYM(3,9) = \TSYM(9,3) = 3,\qquad
  \TSYM(4,9) = \TSYM(9,4) = 7,
\]
where the upper-triangle entries are authoritative under SYM. The
single asymmetric cell $\TRAW(9,4) = 3$ kills the size-$7$ shell
$\{0,4,5,6,7,8,9\}$ under $(\TRAW, B)$, while $(\TSYM, B)$ admits
it. Theorem \ref{thm:lens-internal} below shows that this
phenomenon is invisible at the three-table level: regardless of
the $T$-lens choice, the size-$7$ shell is jointly closed under
$(T,B,C)$ if and only if it is jointly closed under $(\TSYM, B)$,
because $C$ already imposes the corresponding closure
restriction.

\paragraph{Outline.}

Section~\ref{sec:tables} fixes the three tables. Section
\ref{sec:enumeration} performs the brute-force enumeration and
proves Theorem~\ref{thm:three-table-chain}.
Section~\ref{sec:lens-internal} treats the lens-internal
phenomenon at size $7$.
Section~\ref{sec:invariance} records the lens- and
table-invariance of the four-core and the attractor.
Section~\ref{sec:discussion} closes with a brief reflection on
what the three-table chain says about the architecture.

\section{The three tables}
\label{sec:tables}

\begin{definition}[$T$, $B$, $C$]
\label{def:tables}
Let $T,B,C : \Z/10\Z\times\Z/10\Z \to \Z/10\Z$ be the three
$10\times 10$ tables decoded from the canonical bit pattern
\texttt{CL\_BIT\_PATTERN}, with HARMONY-cell counts $73$, $28$,
and $44$ respectively. The tables are recorded in
\cite{Gen13Foundations}; for the present paper we treat them as
fixed input data.
\end{definition}

For the (canonical commutative) version of $T$ we take
$T = \TSYM$, where $T(i,j) = \TRAW(i,j)$ for $i\leq j$ and
$T(j,i) = T(i,j)$ for $i<j$. Section~\ref{sec:lens-internal}
treats the relationship to the non-commutative
$\TRAW$ explicitly.

\begin{lemma}[Individual closed sub-magma counts]
\label{lem:individual-counts}
Among the $1023$ non-empty subsets of $\{0,\dots,9\}$:
\begin{itemize}[leftmargin=*]
\item exactly $449$ are closed under $T$;
\item exactly $9$ are closed under $B$;
\item exactly $50$ are closed under $C$.
\end{itemize}
\end{lemma}

\begin{proof}
Direct enumeration; verified in \texttt{sfm\_q1\_q6\_q7.py}
\cite{SFM2026}, runtime under $2$ seconds.
\end{proof}

\section{The three-substrate joint chain}
\label{sec:enumeration}

\begin{theorem}[Pairwise and triple joint closures]
\label{thm:joint-counts}
Among the $1023$ non-empty subsets of $\{0,\dots,9\}$:
\begin{itemize}[leftmargin=*]
\item exactly $8$ are jointly closed under $(T,B)$;
\item exactly $49$ are jointly closed under $(T,C)$;
\item exactly $9$ are jointly closed under $(B,C)$;
\item exactly $8$ are jointly closed under all of $T$, $B$, $C$.
\end{itemize}
\end{theorem}

\begin{proof}
Direct enumeration via \texttt{sfm\_q1\_q6\_q7.py} \cite{SFM2026}.
The runtime is under $2$ seconds; the script prints each joint
count and the chain at the end.
\end{proof}

The two equalities $|T \cap B| = |T \cap B \cap C| = 8$ are the
load-bearing observation. They say that $C$ does not eliminate
any of the $8$ shells of the $(T,B)$ chain; equivalently, every
$(T,B)$-joint-closed subset is automatically $C$-closed. We
spell this out:

\begin{corollary}[$C$ contains the $(T,B)$ chain]
\label{cor:c-contains-tb}
Every subset jointly closed under $(T,B)$ is closed under $C$.
\end{corollary}

\begin{proof}
The $8$ subsets jointly closed under $(T,B)$ are listed in
Theorem~\ref{thm:three-table-chain}. Each one is verified to be
$C$-closed by direct inspection of $C$ on the corresponding
sub-table; verification in \cite{SFM2026}.
\end{proof}

\begin{corollary}[$C$ extends $T$ on closures]
\label{cor:c-extends-t}
Of the $50$ subsets closed under $C$, exactly $49$ are also closed
under $T$. The unique $C$-closed subset that is not $T$-closed is
recorded in \cite{SFM2026}; it is also not jointly closed under
$(T,B)$ and therefore does not affect Theorem
\ref{thm:three-table-chain}.
\end{corollary}

\begin{proof}
Direct enumeration via \cite{SFM2026}.
\end{proof}

\noindent\emph{Proof of Theorem~\ref{thm:three-table-chain}.}
By Theorem~\ref{thm:joint-counts}, the three-way joint closure
has exactly $8$ elements. By Corollary~\ref{cor:c-contains-tb},
each of the $8$ subsets jointly closed under $(T,B)$ is also
$C$-closed; therefore the three-way joint closure contains all
$8$ subsets of the $(T,B)$ chain. Equality follows by counting.
The chain property (each shell contains the previous) is verified
directly in \cite{SFM2026}. Sizes $\{2,3\}$ are forbidden because
no subset of those sizes is jointly closed under $(T,B)$
\cite{Sanders2026FourCore}, hence none is jointly closed under
$(T,B,C)$. $\square$

\section{Lens-dependence at size 7 is internal to $T$}
\label{sec:lens-internal}

We now relate the three-table picture to the previously-recorded
lens-dependence at size $7$ for the pair $(T,B)$.

\begin{lemma}[The two asymmetric cells of $\TRAW$]
\label{lem:asymmetric}
The cells $(3,9)/(9,3)$ and $(4,9)/(9,4)$ are the only positions
where $\TRAW(i,j)\neq\TRAW(j,i)$. The values are
\[
  \TRAW(3,9)=3,\qquad \TRAW(9,3)=7;\qquad
  \TRAW(4,9)=7,\qquad \TRAW(9,4)=3.
\]
At all other $96$ off-diagonal cells, $\TRAW$ is symmetric. The
upper-triangle authoritative symmetrization $\TSYM$ takes
$\TSYM(3,9)=\TSYM(9,3)=3$ and
$\TSYM(4,9)=\TSYM(9,4)=7$.
\end{lemma}

\begin{proof}
Cell-by-cell verification of the canonical bit-pattern decode;
verified at machine precision in \cite{Gen13Foundations}.
\end{proof}

\begin{theorem}[Size-$7$ lens-dependence under $(T,B)$]
\label{thm:size7-lens}
The subset $S_7 := \{0,4,5,6,7,8,9\}$ satisfies:
\begin{itemize}[leftmargin=*]
\item $S_7$ is jointly closed under $(\TSYM, B)$;
\item $S_7$ is \emph{not} closed under $\TRAW$, because
  $\TRAW(9,4) = 3 \notin S_7$.
\end{itemize}
All other shells of the $8$-element chain (sizes
$\{1,4,5,6,8,9,10\}$) are lens-invariant: closed under both
$\TRAW$ and $\TSYM$.
\end{theorem}

\begin{proof}
Closure of $S_7$ under $\TSYM$ and $B$: cell-by-cell verification
on $S_7\times S_7$, all $49$ cells satisfy
$\TSYM(i,j),B(i,j)\in S_7$.

Failure of $S_7$ under $\TRAW$: the asymmetric cell $(9,4)$
satisfies $9,4\in S_7$ and $\TRAW(9,4)=3\notin S_7$.

Lens-invariance of the other shells: each of $\{0\}$,
$\{0,7,8,9\}$, $\{0,6,7,8,9\}$, $\{0,5,6,7,8,9\}$,
$\{0,3,4,5,6,7,8,9\}$, $\{0,2,3,4,5,6,7,8,9\}$,
$\{0,1,\dots,9\}$ either omits index $3$ or omits index $4$
(or both), so the asymmetric cell $\TRAW(9,4) = 3$ versus
$\TSYM(9,4)=7$ is irrelevant: in each case, either $9\notin S$ or
$4\notin S$ (and both indices $3$ and $4$ are simultaneously
absent for $\{0\}$ and $\{0,7,8,9\}$). The asymmetric cell
$\TRAW(9,3)=7$ versus $\TSYM(9,3)=3$ is similarly contained: it
matters only at the size-$8$ shell $\{0,3,4,5,6,7,8,9\}$, where
both $\TRAW(9,3)=7\in S$ and $\TSYM(9,3)=3\in S$, so neither
violates closure. Verified in \cite{Gen13WP115}.
\end{proof}

\begin{theorem}[Lens-dependence is internal to $T$ at the three-table level]
\label{thm:lens-internal}
Under either $T$-lens choice, the three-substrate joint chain
specializes the corresponding two-substrate joint chain
\emph{verbatim}:
\begin{itemize}[leftmargin=*]
\item under $(\TSYM, B, C)$, the joint chain has $8$ shells of
  sizes $\{1,4,5,6,7,8,9,10\}$ (this is
  Theorem~\ref{thm:three-table-chain});
\item under $(\TRAW, B, C)$, the joint chain has $7$ shells of
  sizes $\{1,4,5,6,8,9,10\}$ — same as under $(\TRAW, B)$.
\end{itemize}
The size-$7$ shell $S_7=\{0,4,5,6,7,8,9\}$ is $C$-closed (in
fact, all $49$ cells of $S_7\times S_7$ map under $C$ into $S_7$,
verified directly), so $C$ does not arbitrate the asymmetric
cell $\TRAW(9,4)=3$: $S_7$ is included in the
$(\TSYM, B, C)$ chain because $\TSYM(9,4)=7\in S_7$, and is
excluded from the $(\TRAW, B, C)$ chain because $\TRAW(9,4)=3
\notin S_7$. The lens-dependence at size $7$ persists at the
three-table level — it is not cured by $C$, nor amplified, nor
shifted to another shell.
\end{theorem}

\begin{proof}
$C$-closure of $S_7$: direct cell-level verification on the
$49$ entries $C(i,j)$ for $(i,j)\in S_7\times S_7$; every value
lies in $S_7=\{0,4,5,6,7,8,9\}$. (The rows of $C$ at indices
$0,4,5,6,7,8,9$ restricted to columns
$0,4,5,6,7,8,9$ map into $S_7$; verification in
\cite{SFM2026}.) Hence $S_7$ is $C$-closed.

Because $S_7$ is $C$-closed but \emph{not} $\TRAW$-closed
(Theorem~\ref{thm:size7-lens}, Lemma~\ref{lem:asymmetric}), the
shell is in the $(\TRAW, B)$ joint chain iff it is in the
$(\TRAW, B, C)$ joint chain — both depend only on whether
$\TRAW$ closes on $S_7$, since $B$ and $C$ both close on it. The
two chains therefore coincide at size $7$ and trivially at every
other shell. The $(\TSYM, B, C)$ count is established in
Theorem~\ref{thm:three-table-chain} ($8$ shells); the
$(\TRAW, B, C)$ count is $7$ shells, lacking exactly $S_7$.

The point is a clarification: the lens-dependence is
\emph{internal to $T$}. It survives at the three-table level
exactly because the third substrate $C$ neither cures nor
extends it. The size-$7$ shell sits inside $C$'s $50$-element
closure structure (Corollary~\ref{cor:c-contains-tb}) regardless
of which $T$-lens is in scope; whether it appears at the joint
level depends only on the $T$-lens.

Direct enumeration confirms the chain shapes:
$\{1,4,5,6,7,8,9,10\}$ under $(\TSYM, B, C)$ specializes to
$\{1,4,5,6,8,9,10\}$ under $(\TRAW, B, C)$, with the size-$7$
shell the unique difference. Verified in \cite{SFM2026}.
\end{proof}

\begin{remark}[The single cell, restated]
\label{rem:single-cell}
The lens-dependence is forced by exactly one cell: the
asymmetric entry $\TRAW(9,4) = 3$ versus the upper-triangle
authoritative value $\TSYM(9,4) = 7$. This is the single most
load-bearing arithmetic fact in the lens-dependence story; the
remaining asymmetric cell $\TRAW(9,3) = 7$ versus
$\TSYM(9,3) = 3$ does not change the chain because no shell of
the $8$-shell chain restricts to a closure problem at that
cell. The substrate is gentle in the sense that one of the two
asymmetric cell pairs is invisible to the chain structure; it is
strict in the sense that the other pair forces a one-shell
deletion at one position.
\end{remark}

\section{Lens- and table-invariance of the four-core and attractor}
\label{sec:invariance}

\begin{theorem}[Four-core is a three-table fixed point]
\label{thm:four-core}
The four-core $\{0,7,8,9\}$ is closed under all three tables
$T$, $B$, $C$. The closure is moreover $T$-lens-invariant: it
holds for both $\TRAW$ and $\TSYM$.
\end{theorem}

\begin{proof}
The four-core does not contain index $3$ or index $4$, so the
two asymmetric cells of $\TRAW$ (Lemma~\ref{lem:asymmetric}) lie
entirely outside $\{0,7,8,9\}^2$. Hence $\TRAW$ and $\TSYM$ agree
on the four-core's product, and the four-core is closed under
both. Closure under $B$ and $C$ is verified directly on the
$16$ cells of $\{0,7,8,9\}^2$ \cite{Gen13Foundations,SFM2026}.
\end{proof}

\begin{theorem}[Attractor is a three-table fixed point]
\label{thm:attractor}
The closed-form $T+B$-mix attractor at mixing weight $\alpha=1/2$
on $\Z/10\Z$, established for $(T,B)$ in
\cite{Sanders2026Attractor,Sanders2026FourCore} as
\[
  (p^*_V, p^*_H, p^*_{\mathrm{Br}}, p^*_R) =
  (0.138,\, 0.540,\, 0.198,\, 0.124),
  \quad H/\mathrm{Br} = 1+\sqrt{3} \text{ exactly},
\]
is supported on the four-core $\{0,7,8,9\}$ and is therefore
unaffected by adjoining $C$ as a third table or by the choice of
$T$-lens.
\end{theorem}

\begin{proof}
The attractor places mass $1$ on the four-core; non-four-core
operators carry zero mass at the fixed point
\cite{Sanders2026Attractor}. The four-core is invariant under
$T$ (both lenses), $B$, and $C$ by Theorem~\ref{thm:four-core}.
Hence the dynamical system, restricted to the four-core support,
is identical under any of the three tables individually or
jointly, and the attractor lifts unchanged.
\end{proof}

\begin{remark}[All non-size-$7$ shells are lens-invariant]
\label{rem:other-shells}
By Theorem~\ref{thm:size7-lens}, only size $7$ exhibits
lens-dependence at the $(T,B)$ level; the other seven shells
are lens-invariant. By Theorem~\ref{thm:lens-internal}, this
remains true at the three-table level. The lens-dependence is
concentrated at exactly one position, and it neither propagates
to other shells nor is cured by the third substrate.
\end{remark}

\section{Discussion: a tiny non-commutativity, located precisely}
\label{sec:discussion}

The structural curiosity recorded in this note is small in
absolute size and sharp in scope. A $10\times 10$ table $T$ with
two asymmetric cell pairs admits two principled symmetrizations,
$\TRAW$ and $\TSYM$. A second $10\times 10$ table $B$ paired
with each of these gives joint chains of $7$ and $8$ shells
respectively, with a single shell of difference at a single size.
A third table $C$, independently constructed and bearing $50$
closed sub-magmas of its own, joins the picture to give a
three-table joint chain of exactly $8$ shells under
$(\TSYM, B, C)$ and $7$ under $(\TRAW, B, C)$ — in both cases
identical to the corresponding two-table chain.

What this says, in one sentence: \emph{the lens-dependence at
size $7$ is internal to $T$, not a property of the joint
substrate}. The $C$ table, despite its $50$-element individual
closure count, sits comfortably above the $8$-element chain that
$T$ and $B$ alone already define; adding $C$ neither sharpens
nor relaxes the lens choice. The non-commutative cell
$\TRAW(9,4)=3$ is the single point of structural difference, and
the third substrate honors that difference without arbitrating
it.

A natural reading is that $C$, $T$, and $B$ are three legitimate
$\Z/10\Z$ multiplications drawn from the same canonical
bit-pattern encoding, and that the chain $\{1,4,5,6,7,8,9,10\}$
is a feature of the encoding itself, not of any single table.
The chain survives the addition of an independently-constructed
third substrate; it survives the lens choice as long as the
lens choice is internalized at the level of $T$. Whether this
chain extends to a fourth or fifth substrate of the same
provenance is open.

\section*{Verification scripts}

The script \texttt{sfm\_q1\_q6\_q7.py} \cite{SFM2026} performs
all the substrate enumerations of this paper at machine
precision: it computes individual closed sub-magma counts under
$T$, $B$, $C$ ($449$, $9$, $50$ respectively); pairwise joint
counts ($|T\cap B|=8$, $|T\cap C|=49$, $|B\cap C|=9$); the
three-way joint count ($|T\cap B\cap C|=8$); and the explicit
$8$-shell chain. Independently, \texttt{joint\_chain\_attractor.py}
\cite{Gen13WP115} reproduces the $(T,B)$ count and chain. Both
runtimes are under $10$ seconds and deterministic.

\begin{thebibliography}{99}

\bibitem{DrapalWanless2021}
A.~Drápal and I.\,M.\ Wanless, \emph{Maximally non-associative
quasigroups}, J.\ Combin.\ Theory Ser.\ A \textbf{184} (2021),
105510.

\bibitem{Sanders2026CLAxioms}
B.\,R.\ Sanders and M.\ Gish, \emph{The CL Forcing Axioms: A1--A9
Uniquely Force the Canonical Composition Lattice on $\Z/10\Z$},
in preparation, 2026.

\bibitem{Sanders2026FourCore}
B.\,R.\ Sanders and M.\ Gish, \emph{Joint Closure, Per-Coordinate
Fuse Data, and a Closed-Form Algebraic Attractor of Two
Commutative Binary Operations on $\Z/10\Z$}, submitted to
\emph{Algebraic Combinatorics}, 2026 (J02 in the release
sequence).

\bibitem{Sanders2026LensInvariance}
B.\,R.\ Sanders and M.\ Gish, \emph{TSML 73 Cells / BHML 28
Cells: Lens-Invariant Cell Counts on the $\Z/10\Z$ Composition
Lattice}, submitted to \emph{Experimental Mathematics}, 2026
(J05 in the release sequence).

\bibitem{Sanders2026Attractor}
B.\,R.\ Sanders and M.\ Gish, \emph{Closed-Form Attractor and
the $\alpha=1/2$ Algebraic Singularity}, in preparation, 2026.

\bibitem{Gen13WP115}
B.\,R.\ Sanders et al., \emph{WP115: The Joint TSML+BHML
Closed-Subset Chain and the Universal Four-Core Attractor —
Verification Script \texttt{joint\_chain\_attractor.py}},
codebase release 2026.

\bibitem{Gen13Foundations}
B.\,R.\ Sanders et al., \emph{Gen13 Foundations Module:
\texttt{cl.py} with \texttt{CL\_BIT\_PATTERN},
\texttt{CL\_TSML\_RAW}, \texttt{CL\_TSML\_SYM},
\texttt{CL\_BHML}, \texttt{CL\_STD}, \texttt{lenses.py}},
codebase release 2026.

\bibitem{SFM2026}
B.\,R.\ Sanders, M.\ Gish, et al., \emph{Substrate Function Map:
Q1--Q6--Q7 Verification (\texttt{sfm\_q1\_q6\_q7.py})},
\texttt{Atlas/META\_PLAN\_2026-05-06/SUBSTRATE\_FUNCTION\_MAP/},
codebase release 2026-05-08.

\end{thebibliography}

\end{document}
